\section*{الفصل \ref{ch:g12:satellites-kepler} --- التوابع وحركة الكواكب}

\begin{solution}{exo:g12:satellites-kepler:1}
على المنصة: $F = GMm/R^2 = \num{3.98e17}/\num{4.06e13} \approx
\qty{9.8e3}{N}$. وعلى $r = \qty{6.77e6}{m}$:
$F \approx \qty{8.7e3}{N}$ --- أي ما يزال $89\%$ من القيمة الأرضية. ويطفو
من على المتن لأنهم \emph{يسقطون} مع المحطة،
لا لأن \omterm{def:g11:fundamental-interactions:four}{الجاذبية} زالت.
\end{solution}

\begin{solution}{exo:g12:satellites-kepler:2}
$v = \sqrt{GM/R} = \sqrt{\num{3.98e14}/\num{6.37e6}} \approx
\qty{7.9e3}{m/s}$؛ $T = 2\pi R/v \approx \qty{5.1e3}{s} \approx
\qty{84}{min}$. وهي بالضبط \qty{7.9}{km/s} لكرة المدفع:
\omterm{def:g10:universal-gravitation:satellite}{فالمدار} الملامس \emph{هو} رمية نيوتن.
\end{solution}

\begin{solution}{exo:g12:satellites-kepler:3}
$v \propto 1/\sqrt{r}$: أي مقسومة على $\sqrt 2 \approx 1.41$.
و$T \propto r^{3/2}$: أي مضروب في $2^{3/2} \approx 2.8$. ولا يحوي أيّ منهما
$m$: فلا.
\end{solution}

\begin{solution}{exo:g12:satellites-kepler:4}
$v = \sqrt{\num{3.98e14}/\num{6.77e6}} \approx \qty{7.67e3}{m/s}$؛
$T = 2\pi r/v \approx \qty{5.55e3}{s} \approx \qty{92}{min}$؛
و$\num{86400}/5546 \approx 15.6$ مدارًا في اليوم.
\end{solution}

\begin{solution}{exo:g12:satellites-kepler:5}
إهليلجات والشمس في إحدى البؤرتين؛ ومساحات متساوية في أزمنة متساوية؛
و$T^2/a^3$ متطابقة لكل الكواكب. وأسرع ما يكون عند \omterm{def:g12:satellites-kepler:ellipse}{الحضيض}، بالقانون
الثاني (فقطعة الكنس القصيرة يجب أن تكنس سريعًا). ولا: فكتلة \omterm{def:g10:universal-gravitation:satellite}{التابع}
تتلاشى --- ولا تدخل إلا الكتلة المركزية.
\end{solution}

\begin{solution}{exo:g12:satellites-kepler:6}
$T = \qty{2.36e6}{s}$: $M = 4\pi^2 r^3/(GT^2) =
\num{2.24e27}/371 \approx \qty{6.0e24}{kg}$ --- أي في حدود $1\%$ من
\qty{5.97e24}{kg}.
\end{solution}

\begin{solution}{exo:g12:satellites-kepler:7}
تتجه \omterm{def:g11:fundamental-interactions:four}{الجاذبية} نحو مركز الأرض، وفي الحركة الدائرية تتجه \omterm{def:g11:forces-and-motion:netforce}{القوة
المحصّلة} نحو مركز الدائرة: فينطبق المركزان. أما
دائرة فوق دائرة عرض باريس فمركزها على المحور، لا مركز
الأرض --- وذلك مستحيل. ولا يفلح التعلّق إلا حيث تكون دائرة العرض
دائرة عظمى مركزها المركز: أي فوق \emph{خط الاستواء}.
\end{solution}

\begin{solution}{exo:g12:satellites-kepler:8}
$T = \qty{2.754e4}{s}$: $M = 4\pi^2 r^3/(GT^2) =
\num{3.26e22}/\num{5.06e-2} \approx \qty{6.4e23}{kg}$ --- أي مريخ بطاقة
المعطيات.
\end{solution}

\begin{solution}{exo:g12:satellites-kepler:9}
الأرض: $T^2/a^3 = \num{9.96e14}/\num{3.35e33} =
\qty{2.98e-19}{s^2/m^3}$. والمريخ: $T = \qty{5.94e7}{s}$،
$\num{3.52e15}/\num{1.18e34} = \qty{2.98e-19}{s^2/m^3}$ --- وهما متساويان، كما
يقتضي القانون الثالث. ومنه $M = 4\pi^2/(G \times \num{2.98e-19})
\approx \qty{1.99e30}{kg}$.
\end{solution}

\begin{solution}{exo:g12:satellites-kepler:10}
$v = 2\pi r/T = \qty{1.02e3}{m/s}$؛ $a = v^2/r =
\qty{2.72e-3}{m/s^2}$؛ $g/60^2 = 9.81/3600 = \qty{2.73e-3}{m/s^2}$.
فالقمر يقع على بعد $60$ نصف قطر أرضي، وسقوطه أضعف $60^2$ مرة
من سقوط التفاحة: أي بمربع المسافة في المقام بالضبط --- فقانون واحد
لكليهما.
\end{solution}

\begin{solution}{exo:g12:satellites-kepler:11}
سيريس: $T = 2.77^{3/2} = \sqrt{21.3} \approx \qty{4.6}{yr}$. وهالي:
$T = 17.8^{3/2} \approx \qty{75}{yr}$، ومنه يكون \omterm{def:g12:satellites-kepler:ellipse}{الحضيض} التالي نحو
$1986 + 75 \approx 2061$.
\end{solution}

\begin{solution}{exo:g12:satellites-kepler:12}
عند \qty{400}{km}: $v \approx \qty{7.67e3}{m/s}$؛ وعند \qty{350}{km}
($r = \qty{6.72e6}{m}$): $v \approx \qty{7.70e3}{m/s}$. \omterm{def:g10:inertia:inventory}{والاحتكاك} يستنزف
\omterm{def:g11:mechanical-energy:em}{الطاقة الميكانيكية} فعلًا --- فيتقلص \omterm{def:g10:universal-gravitation:satellite}{المدار}؛ لكن في طريق النزول
يردّ شغل \omterm{def:g11:fundamental-interactions:four}{الجاذبية} أكثر مما تخسره \omterm{def:g12:kinematics-2d:velocity}{السرعة} \omterm{def:g11:circuits-and-power:resistance}{بالمقاومة}:
فيهبط $E_m$ بينما ترتفع $E_k$.
\end{solution}

\begin{solution}{exo:g12:satellites-kepler:13}
$T = \qty{3.024e5}{s}$: أي $M = 4\pi^2 r^3/(GT^2) =
\num{1.65e31}/6.10 \approx \qty{2.7e30}{kg} \approx 1.4$ شمسًا. أما كتلة
\emph{الكوكب} فبعيدة المنال: إذ تلاشت من معادلات
\omterm{def:g10:universal-gravitation:satellite}{المدار}.
\end{solution}

\begin{solution}{exo:g12:satellites-kepler:14}
$T = \qty{8.862e4}{s}$، $GM = \qty{4.28e13}{m^3/s^2}$:
$r^3 = GMT^2/(4\pi^2) = \num{8.52e21}$، ومنه $r \approx
\qty{2.04e7}{m}$ و$h = r - R \approx \qty{1.7e7}{m} \approx
\qty{17000}{km}$ فوق سطح المريخ.
\end{solution}

\begin{solution}{exo:g12:satellites-kepler:15}
$v_a = r_p v_p / r_a = \num{8.77e10} \times \num{5.45e4} /
\num{5.25e12} \approx \qty{9.1e2}{m/s}$. وبما أن $r_a/r_p \approx 60$،
يكون المذنّب أبعد $60$ مرة وأبطأ $60$ مرة: فيقضي
عقودًا يتباطأ في الظلام وأسابيع ينطلق أمام الشمس.
\end{solution}

\begin{solution}{pb:g12:satellites-kepler:1}
\textbf{1.} $T = \qty{86164}{s}$: فالأرض تدور
\emph{بالنسبة إلى النجوم} دورةً في \omterm{def:g12:satellites-kepler:geo}{اليوم النجمي}؛ أما اليوم \qty{24}{h}
فيضيف الدرجة في اليوم التي تتقدمها الأرض حول الشمس.

\textbf{2.} \omterm{def:g11:forces-and-motion:netforce}{القوة المحصّلة} في الحركة الدائرية مصوَّبة نحو مركز
الدائرة؛ \omterm{def:g10:inertia:force}{والقوة} الوحيدة، وهي \omterm{def:g11:fundamental-interactions:four}{الجاذبية}، مصوَّبة نحو مركز الأرض: فالمركزان
نقطة واحدة.

\textbf{3.} ودائرة تتعلّق فوق باريس سيكون مركزها على المحور
شمال المركز --- وهو مستبعَد بالسؤال 2. فيجب أن يقع \omterm{def:g10:universal-gravitation:satellite}{المدار} في
\emph{المستوي الاستوائي}.

\textbf{4.} المركّبة الناظمية \omterm{thm:g12:newtons-laws:second}{لقانون نيوتن الثاني}:
$m v^2/r = GmM/r^2$، ومنه $v = \sqrt{GM/r}$ (فتتلاشى $m$).

\textbf{5.} $T = 2\pi r/v = 2\pi\sqrt{r^3/(GM)}$؛ وبالتربيع،
$T^2/r^3 = 4\pi^2/(GM)$.

\textbf{6.} $r = \left(GMT^2/4\pi^2\right)^{1/3} =
(\num{7.49e22})^{1/3} \approx \qty{4.22e7}{m}$.

\textbf{7.} $h = r - R = \num{4.22e7} - \num{6.37e6} \approx
\qty{3.58e7}{m} \approx \qty{35800}{km}$.

\textbf{8.} $v = 2\pi \times \num{4.22e7}/\num{86164} \approx
\qty{3.07e3}{m/s}$.

\textbf{9.} $\sqrt{\num{3.98e14}/\num{4.22e7}} \approx
\qty{3.07e3}{m/s}$ --- وهو متسق.

\textbf{10.} نصفا القطر متطابقان: فالكتلة تلاشت في السؤال 4. \omterm{def:g10:universal-gravitation:satellite}{والتابع}
الأثقل لا يغيّر إلا فاتورة الإطلاق --- أي مزيدًا من الوقود \omterm{def:g10:universal-gravitation:satellite}{للمدار} نفسه.

\textbf{11.} $v = \sqrt{\num{3.98e14}/\num{6.77e6}} \approx
\qty{7.67e3}{m/s}$.

\textbf{12.} $T = 2\pi r/v \approx \qty{5.55e3}{s} \approx
\qty{92.4}{min}$.

\textbf{13.} $\num{86400}/5546 \approx 15.6$ مدارًا: أي نحو $16$
شروقًا في اليوم.

\textbf{14.} $r_{\text{الثابت}}/r_{\text{المحطة}} = \num{4.22e7}/\num{6.77e6}
= 6.23$؛ و$6.23^{3/2} \approx 15.5 = \num{86164}/5546$ --- أي قانون كبلر
الثالث، متحقَّقًا منه في الحال.

\textbf{15.} المحطة أسرع (\qty{7.67}{km/s} مقابل
\qty{3.07}{km/s})، ومع ذلك كلّف الثابت بالنسبة إلى الأرض \omterm{def:g10:energy-conservation:energy}{طاقة} أكبر لكل
\omterm{def:g10:orders-of-magnitude:unit}{كيلوغرام}: فالصعود في \omterm{def:g10:energy-conservation:energy}{الطاقة} الكامنة يفوق خسارة \omterm{def:g12:kinematics-2d:velocity}{السرعة}.

\textbf{16.} $M_J = 4\pi^2 r^3/(G T^2)$.

\textbf{17.} $T = \qty{1.529e5}{s}$: $M_J =
\num{2.97e27}/1.56 \approx \qty{1.90e27}{kg}$.

\textbf{18.} $\num{1.90e27}/\num{5.97e24} \approx 318$ أرضًا ---
ونحو $1/1000$ من الشمس.

\textbf{19.} $T_J = \qty{3.576e4}{s}$: $r^3 = GM_J T_J^2/(4\pi^2) =
\num{4.11e24}$، $r \approx \qty{1.60e8}{m}$ --- أما آيو، على
\qty{4.22e8}{m}، فيطير أعلى $2.6$ مرة وينجرف عبر سماء
المشتري.

\textbf{20.} اركن المرحّل على الدائرة الاستوائية ذات الارتفاع
\qty{35800}{km}، حيث يركب \omterm{def:g12:kinematics-2d:velocity}{بسرعة} \qty{3.07}{km/s} ولا يغادر
تصويب الطبق أبدًا؛ والقانون نفسه، إذا غُذّي بشهر آيو، \omterm{def:g10:universal-gravitation:weight}{وزن} المشتري عند
$\qty{1.90e27}{kg}$ --- أي $318$ أرضًا --- مجانًا.
\end{solution}
